\documentclass[12pt,a4paper]{article}

% ---- Encoding & Font ----
\usepackage{iftex}
\ifPDFTeX
  \usepackage[utf8]{inputenc}
  \usepackage[T1]{fontenc}
\else
  \usepackage{fontspec}  % XeTeX/LuaTeX (tectonic): Unicode nativo (·, —, ¿, ª, §)
\fi
\usepackage[spanish]{babel}
\usepackage{csquotes}

% ---- Page layout ----
\usepackage[top=2.5cm, bottom=2.5cm, left=2.5cm, right=2.5cm]{geometry}
\usepackage{setspace}
\onehalfspacing

% ---- Math ----
\usepackage{amsmath, amssymb, amsthm}
\usepackage{mathtools}
\usepackage{physics}

% ---- Graphics ----
\usepackage{graphicx}
\usepackage{tikz}
\usepackage{float}
\usepackage{caption}

% ---- Tables ----
\usepackage{array}
\usepackage{booktabs}
\usepackage{tabularx}
\usepackage{colortbl}

% ---- Colours ----
\usepackage{xcolor}
\definecolor{notebg}{HTML}{FFF3CD}
\definecolor{noteborder}{HTML}{FFC107}
\definecolor{boxred}{HTML}{F8D7DA}
\definecolor{boxredborder}{HTML}{F5C6CB}

% ---- Boxes ----
\usepackage{mdframed}
\usepackage{tcolorbox}
\tcbuselibrary{most}

\newtcolorbox{keybox}[1][]{
  colback=blue!5,
  colframe=blue!70!black,
  arc=4pt,
  boxrule=1pt,
  fonttitle=\bfseries,
  title={#1}
}

\newtcolorbox{warnbox}[1][]{
  colback=boxred,
  colframe=boxredborder,
  coltitle=black!75,
  arc=4pt,
  boxrule=1pt,
  fonttitle=\bfseries,
  title={#1}
}

\newtcolorbox{notebox}[1][]{
  colback=notebg,
  colframe=noteborder,
  coltitle=black!75,
  arc=4pt,
  boxrule=1pt,
  fonttitle=\bfseries,
  title={#1}
}

% ---- Diagramas (gráficas) ----
\usepackage{pgfplots}
\pgfplotsset{compat=1.18}
\usetikzlibrary{babel}  % babel español activa `>`; esto evita romper las flechas `->`
% courses/CDIV2017/Clases/assets/tikzstyles.tex
% ---------------------------------------------------------------------------
% Estilos compartidos para los diagramas del curso Cálculo 1 (CDIV2017).
% Fuente única del "look" visual (colores, grosores, escalas) para que todas
% las figuras (TikZ / pgfplots: conjuntos, rectas numéricas, gráficas de
% funciones) sean consistentes entre clases.
%
% Uso: la clase debe cargar antes `tikz` y `pgfplots` (bloque §6.3 del
%      CLAUDE.md del curso), y luego
%      % courses/CDIV2017/Clases/assets/tikzstyles.tex
% ---------------------------------------------------------------------------
% Estilos compartidos para los diagramas del curso Cálculo 1 (CDIV2017).
% Fuente única del "look" visual (colores, grosores, escalas) para que todas
% las figuras (TikZ / pgfplots: conjuntos, rectas numéricas, gráficas de
% funciones) sean consistentes entre clases.
%
% Uso: la clase debe cargar antes `tikz` y `pgfplots` (bloque §6.3 del
%      CLAUDE.md del curso), y luego
%      % courses/CDIV2017/Clases/assets/tikzstyles.tex
% ---------------------------------------------------------------------------
% Estilos compartidos para los diagramas del curso Cálculo 1 (CDIV2017).
% Fuente única del "look" visual (colores, grosores, escalas) para que todas
% las figuras (TikZ / pgfplots: conjuntos, rectas numéricas, gráficas de
% funciones) sean consistentes entre clases.
%
% Uso: la clase debe cargar antes `tikz` y `pgfplots` (bloque §6.3 del
%      CLAUDE.md del curso), y luego
%      \input{../assets/tikzstyles.tex}
% (compilar desde el directorio de la clase para que la ruta relativa resuelva).
%
% Paleta propia de este curso (no se reusa la de courses/Fisica3-2015): mismos
% tonos base (azul/rojo/ámbar) para alinear con keybox/notebox/warnbox, pero
% archivo independiente, sin dependencia cruzada entre cursos.
% ---------------------------------------------------------------------------

% ---- Paleta (alineada con las cajas del preámbulo: keybox/notebox/warnbox) ----
\definecolor{figblue}{HTML}{1F4E79}   % azul principal (líneas, keybox, conjuntos)
\definecolor{figred}{HTML}{B03A2E}    % rojo de énfasis (warnbox, región resaltada)
\definecolor{figamber}{HTML}{B8860B}  % ámbar (notebox)
\definecolor{figgray}{HTML}{555555}   % auxiliares, ejes, guías, universo

% ---- TikZ: librerías y estilos reutilizables ----
% Nota: las puntas se declaran como `-{Latex[...]}` y no como `->`. La forma
% con llaves NO contiene un `>` literal, así que es inmune al `>` activo que
% introduce babel español (de todos modos se carga \usetikzlibrary{babel} en
% el bloque §6.3 para cubrir cualquier `->` suelto).
\usetikzlibrary{arrows.meta,calc,angles,quotes,patterns}

\tikzset{
  figline/.style={figblue, line width=0.9pt},
  figaux/.style={figgray, dashed, line width=0.6pt},
  figaxis/.style={figgray, line width=0.7pt, -{Latex[length=2.2mm,width=1.5mm]}},
  % conjuntos (diagramas de Venn / patatoidales)
  figset/.style={figline, fill=figblue!6},
  figsetB/.style={figred, line width=0.9pt, fill=figred!6},
  figsetC/.style={figamber, line width=0.9pt, fill=figamber!10},
  figuniverso/.style={figgray, line width=0.7pt},
  % puntos de la recta / conjuntos numéricos
  figpto/.style={circle, inner sep=0pt, minimum size=5pt, draw=none},
  figptoabierto/.style={circle, inner sep=0pt, minimum size=5pt, draw=figblue, line width=0.9pt, fill=white},
  figptocerrado/.style={circle, inner sep=0pt, minimum size=5pt, draw=figblue, line width=0.9pt, fill=figblue},
  figptoY/.style={figpto, fill=figblue},
  figptoZ/.style={figpto, fill=figred},
  % segmentos gruesos para intervalos sobre la recta
  figintervalo/.style={figblue, line width=2.2pt},
  % vectores/flechas de transformación (traslación, dilatación)
  figvec/.style={figred, line width=1pt, -{Latex[length=2.2mm,width=1.6mm]}},
  % etiquetas
  figlbl/.style={font=\small},
  figlblaux/.style={font=\footnotesize, figgray},
}

% ---- pgfplots: estilo base de las gráficas de funciones ----
\pgfplotsset{
  calcfig/.style={
    width=10cm, height=6cm,
    axis lines=middle,
    xlabel style={font=\small, at={(axis description cs:1,0.5)}, anchor=west},
    ylabel style={font=\small, at={(axis description cs:0.5,1)}, anchor=south},
    tick label style={font=\footnotesize},
    every axis plot/.append style={line width=1.1pt, figblue},
    grid=none,
    legend cell align=left,
    legend style={font=\footnotesize, draw=figgray!40},
    clip=false,
  },
}

% (compilar desde el directorio de la clase para que la ruta relativa resuelva).
%
% Paleta propia de este curso (no se reusa la de courses/Fisica3-2015): mismos
% tonos base (azul/rojo/ámbar) para alinear con keybox/notebox/warnbox, pero
% archivo independiente, sin dependencia cruzada entre cursos.
% ---------------------------------------------------------------------------

% ---- Paleta (alineada con las cajas del preámbulo: keybox/notebox/warnbox) ----
\definecolor{figblue}{HTML}{1F4E79}   % azul principal (líneas, keybox, conjuntos)
\definecolor{figred}{HTML}{B03A2E}    % rojo de énfasis (warnbox, región resaltada)
\definecolor{figamber}{HTML}{B8860B}  % ámbar (notebox)
\definecolor{figgray}{HTML}{555555}   % auxiliares, ejes, guías, universo

% ---- TikZ: librerías y estilos reutilizables ----
% Nota: las puntas se declaran como `-{Latex[...]}` y no como `->`. La forma
% con llaves NO contiene un `>` literal, así que es inmune al `>` activo que
% introduce babel español (de todos modos se carga \usetikzlibrary{babel} en
% el bloque §6.3 para cubrir cualquier `->` suelto).
\usetikzlibrary{arrows.meta,calc,angles,quotes,patterns}

\tikzset{
  figline/.style={figblue, line width=0.9pt},
  figaux/.style={figgray, dashed, line width=0.6pt},
  figaxis/.style={figgray, line width=0.7pt, -{Latex[length=2.2mm,width=1.5mm]}},
  % conjuntos (diagramas de Venn / patatoidales)
  figset/.style={figline, fill=figblue!6},
  figsetB/.style={figred, line width=0.9pt, fill=figred!6},
  figsetC/.style={figamber, line width=0.9pt, fill=figamber!10},
  figuniverso/.style={figgray, line width=0.7pt},
  % puntos de la recta / conjuntos numéricos
  figpto/.style={circle, inner sep=0pt, minimum size=5pt, draw=none},
  figptoabierto/.style={circle, inner sep=0pt, minimum size=5pt, draw=figblue, line width=0.9pt, fill=white},
  figptocerrado/.style={circle, inner sep=0pt, minimum size=5pt, draw=figblue, line width=0.9pt, fill=figblue},
  figptoY/.style={figpto, fill=figblue},
  figptoZ/.style={figpto, fill=figred},
  % segmentos gruesos para intervalos sobre la recta
  figintervalo/.style={figblue, line width=2.2pt},
  % vectores/flechas de transformación (traslación, dilatación)
  figvec/.style={figred, line width=1pt, -{Latex[length=2.2mm,width=1.6mm]}},
  % etiquetas
  figlbl/.style={font=\small},
  figlblaux/.style={font=\footnotesize, figgray},
}

% ---- pgfplots: estilo base de las gráficas de funciones ----
\pgfplotsset{
  calcfig/.style={
    width=10cm, height=6cm,
    axis lines=middle,
    xlabel style={font=\small, at={(axis description cs:1,0.5)}, anchor=west},
    ylabel style={font=\small, at={(axis description cs:0.5,1)}, anchor=south},
    tick label style={font=\footnotesize},
    every axis plot/.append style={line width=1.1pt, figblue},
    grid=none,
    legend cell align=left,
    legend style={font=\footnotesize, draw=figgray!40},
    clip=false,
  },
}

% (compilar desde el directorio de la clase para que la ruta relativa resuelva).
%
% Paleta propia de este curso (no se reusa la de courses/Fisica3-2015): mismos
% tonos base (azul/rojo/ámbar) para alinear con keybox/notebox/warnbox, pero
% archivo independiente, sin dependencia cruzada entre cursos.
% ---------------------------------------------------------------------------

% ---- Paleta (alineada con las cajas del preámbulo: keybox/notebox/warnbox) ----
\definecolor{figblue}{HTML}{1F4E79}   % azul principal (líneas, keybox, conjuntos)
\definecolor{figred}{HTML}{B03A2E}    % rojo de énfasis (warnbox, región resaltada)
\definecolor{figamber}{HTML}{B8860B}  % ámbar (notebox)
\definecolor{figgray}{HTML}{555555}   % auxiliares, ejes, guías, universo

% ---- TikZ: librerías y estilos reutilizables ----
% Nota: las puntas se declaran como `-{Latex[...]}` y no como `->`. La forma
% con llaves NO contiene un `>` literal, así que es inmune al `>` activo que
% introduce babel español (de todos modos se carga \usetikzlibrary{babel} en
% el bloque §6.3 para cubrir cualquier `->` suelto).
\usetikzlibrary{arrows.meta,calc,angles,quotes,patterns}

\tikzset{
  figline/.style={figblue, line width=0.9pt},
  figaux/.style={figgray, dashed, line width=0.6pt},
  figaxis/.style={figgray, line width=0.7pt, -{Latex[length=2.2mm,width=1.5mm]}},
  % conjuntos (diagramas de Venn / patatoidales)
  figset/.style={figline, fill=figblue!6},
  figsetB/.style={figred, line width=0.9pt, fill=figred!6},
  figsetC/.style={figamber, line width=0.9pt, fill=figamber!10},
  figuniverso/.style={figgray, line width=0.7pt},
  % puntos de la recta / conjuntos numéricos
  figpto/.style={circle, inner sep=0pt, minimum size=5pt, draw=none},
  figptoabierto/.style={circle, inner sep=0pt, minimum size=5pt, draw=figblue, line width=0.9pt, fill=white},
  figptocerrado/.style={circle, inner sep=0pt, minimum size=5pt, draw=figblue, line width=0.9pt, fill=figblue},
  figptoY/.style={figpto, fill=figblue},
  figptoZ/.style={figpto, fill=figred},
  % segmentos gruesos para intervalos sobre la recta
  figintervalo/.style={figblue, line width=2.2pt},
  % vectores/flechas de transformación (traslación, dilatación)
  figvec/.style={figred, line width=1pt, -{Latex[length=2.2mm,width=1.6mm]}},
  % etiquetas
  figlbl/.style={font=\small},
  figlblaux/.style={font=\footnotesize, figgray},
}

% ---- pgfplots: estilo base de las gráficas de funciones ----
\pgfplotsset{
  calcfig/.style={
    width=10cm, height=6cm,
    axis lines=middle,
    xlabel style={font=\small, at={(axis description cs:1,0.5)}, anchor=west},
    ylabel style={font=\small, at={(axis description cs:0.5,1)}, anchor=south},
    tick label style={font=\footnotesize},
    every axis plot/.append style={line width=1.1pt, figblue},
    grid=none,
    legend cell align=left,
    legend style={font=\footnotesize, draw=figgray!40},
    clip=false,
  },
}


% ---- Hyperlinks & TOC ----
\usepackage[hidelinks]{hyperref}
\usepackage{bookmark}

% ---- Enumerate ----
\usepackage{enumitem}

% ---- Miscellaneous ----
\usepackage{icomma}
\usepackage{siunitx}
\usepackage{textcomp}
\usepackage[bottom]{footmisc}

% ---- Title ----
\title{\textbf{Clase 1: Presentación del curso y nuevo programa} \\[0.3em]
        \large{Zoología de objetos matemáticos, conjuntos: pertenencia, igualdad e inclusión}}
\author{Transcripto y expandido --- Curso Cálculo 1 (OpenFING)}
\date{}

\begin{document}

\maketitle
\thispagestyle{empty}
\newpage

\tableofcontents
\newpage

\section{Presentación y organización del curso}

\subsection{Docentes y anuncios prácticos}

El docente responsable del teórico de la tarde es \textbf{Alexandre Miquel};
el coordinador del curso y responsable del teórico de la mañana es Marcos
Barrios. El curso se llama oficialmente \textbf{Cálculo DIR} (Cálculo
Diferencial e Integral con una Variable), nombre que corresponde a un
\textbf{nuevo programa} que el docente presenta en esta primera clase.

Anuncios organizativos de la clase:

\begin{itemize}
  \item Registrarse en el sitio \textbf{Eva} del curso.
  \item Hay \textbf{seis grupos de prácticos}; como no existe teórico
    nocturno, el práctico del grupo 1 (nocturno) dura dos horas y media en
    lugar de una hora y media, y en él se hacen repasos de teórico para
    quienes no pueden asistir de día.
  \item Los prácticos llevan, en general, \textbf{una semana de retraso}
    respecto al teórico (es una edición nueva del curso y se están
    ajustando los tiempos).
  \item Habrá un \textbf{cuestionario semanal en línea}, sin valor para la
    nota (no puntúa de forma directa), pensado como termómetro de
    aprendizaje tanto para los estudiantes como para los docentes.
  \item Todavía no hay apuntes nuevos completos: para los temas que no
    cambian respecto al programa anterior se puede seguir usando el
    material viejo; para los temas que sí cambian, el docente producirá
    capítulos específicos.
\end{itemize}

\begin{notebox}[Régimen de exámenes: aún sin definir]
Preguntas de la clase sobre el régimen de exámenes (si las demostraciones
``antiguas'' siguen siendo válidas, si exonerar habilita a dar el examen de
Cálculo 2 viejo, si habrá más parciales que teóricos, etc.) quedaron sin
resolver: el docente respondió explícitamente ``no sé'' en varios casos,
porque es la primera edición del programa nuevo y los detalles de examen
todavía se están coordinando entre los docentes. Un dato que sí quedó
fijado: quien apruebe Cálculo 1 nuevo queda habilitado para dar el examen de
Cálculo 2 nuevo, no el de Cálculo 2 viejo.
\end{notebox}

\subsection{El nuevo programa: ``Cálculo 1 nuevo''}

La reorganización responde a un cambio en los programas de liceo, que dejó
un salto demasiado grande entre el liceo y el temario tradicional de
Cálculo 1. La solución fue \textbf{aligerar Cálculo 1} trasladando temas a
Cálculo 2 y superiores. Importante: los temas no desaparecen del plan de
estudios, solo se \textbf{reubican}.

\textbf{Temas que se van de Cálculo 1} (pasan a Cálculo 2 o cursos
posteriores):

\begin{itemize}
  \item Números complejos.
  \item Sucesiones y series.
  \item Integrales impropias.
\end{itemize}

El nuevo temario se organiza en \textbf{ocho capítulos}:

\begin{table}[H]
\centering
\begin{tabular}{@{}clp{7cm}@{}}
\toprule
\# & Capítulo & Observación \\
\midrule
1 & Conjuntos y funciones & Repaso de notaciones (esta clase) \\
2 & Números reales & Sin cambios respecto al programa anterior; dos semanas \\
3 & \textbf{Integrales} & Tema central de la nueva edición; cambia mucho \\
4 & Continuidad & Cambia mucho (nuevas demostraciones) \\
5 & Funciones trigonométricas & Sin cambios; se dicta justo antes del primer parcial \\
6 & Derivadas & No cambia demasiado \\
7 & Desarrollo de Taylor & --- \\
8 & Métodos de integración & --- \\
\bottomrule
\end{tabular}
\end{table}

El cambio conceptual más importante es que las \textbf{integrales pasan a
ser el tema central} del curso desde el capítulo 3, en vez de dictarse
recién al final (donde antes solía faltar tiempo para cubrirlas bien). Esto
obliga a repensar la presentación, porque en el capítulo 3 todavía no se
dispone ni de sucesiones, ni de continuidad, ni de derivada:

\begin{itemize}
  \item La integral se presenta a partir de la \textbf{teoría de Riemann
    usando solo supremos e ínfimos} (noción que se construye en el
    capítulo 2, números reales), sin necesitar el teorema fundamental del
    cálculo ni la noción de función continua.
  \item Se introduce un teorema nuevo en este enfoque: \textbf{toda función
    monótona es integrable}, que permite construir ejemplos de integrales
    sin pasar por la continuidad.
  \item La integral se entiende, en esta etapa, como el \textbf{área
    algebraica} de la superficie definida por una curva.
  \item El primer parcial cubre normalmente los capítulos 1 a 4.
\end{itemize}

\begin{sloppypar}
En continuidad también hay cambios profundos: los grandes teoremas
clásicos (\textbf{Bolzano}, \textbf{Weierstrass}) se demostraban
tradicionalmente apoyándose en sucesiones. Pero como las sucesiones se
retiraron del programa, hace falta producir \textbf{demostraciones
nuevas que no usen sucesiones}.
\end{sloppypar}

Otra simplificación deliberada: \textbf{no habrá topología de la recta}. El
curso solo considera funciones definidas en intervalos, así que la única
noción topológica necesaria es distinguir intervalo abierto, cerrado,
semiabierto y semicerrado.

Por la centralidad de las integrales, el tema se reparte en varios
capítulos en vez de concentrarse en uno solo: aparece en el capítulo 3
(teoría básica), otra vez en el capítulo 4 (para mostrar que toda función
continua es localmente integrable) y de nuevo en el capítulo 6, derivadas,
donde recién ahí se presenta el \textbf{teorema fundamental del cálculo}.

\section{Zoología de los objetos matemáticos}

Antes de definir conjuntos, el docente hace un repaso de los \textbf{tipos
de objetos matemáticos} que se van a manipular en el curso, usando la
imagen de una ``zoología'' de objetos:

\begin{itemize}
  \item \textbf{Números}: los conjuntos numéricos que se usarán son los
    \textbf{naturales} ($\mathbb{N}$), \textbf{enteros} ($\mathbb{Z}$),
    \textbf{racionales} ($\mathbb{Q}$) y \textbf{reales} ($\mathbb{R}$);
    los \textbf{complejos quedan fuera del programa}. Ejemplos de números:
    $0, 1, 2, -1, 3, \tfrac{1}{2}, \sqrt{2}, \pi, e$.
  \item \textbf{Objetos geométricos}: puntos, rectas, triángulos,
    círculos. En este curso se usan casi exclusivamente como
    \textbf{herramienta para dibujar} funciones, no como objeto de estudio
    en sí.
  \item \textbf{Objetos compuestos}: se construyen agrupando objetos de
    base. Ejemplos: \textbf{funciones} (reglas que transforman un número
    en otro número), matrices, \textbf{pares} $(a,b)$, y \textbf{n-uplas}
    (pares generalizados con cualquier cantidad de componentes).
  \item \textbf{Conjuntos}: la noción central de esta clase, presentada al
    final ``porque es la mejor para el final''.
\end{itemize}

\begin{notebox}[Analogía informática]
La mejor intuición de un conjunto (al menos uno finito) es una
\textbf{carpeta}: contiene ``archivos'' (los elementos). Hay dos
diferencias clave con las carpetas de una computadora: (1) las carpetas
tienen nombre, los \textbf{conjuntos son anónimos} (se identifican
únicamente por lo que contienen); (2) las carpetas son necesariamente
finitas por limitaciones de hardware, mientras que \textbf{los conjuntos
pueden ser infinitos}. El docente aclara que esta cercanía no es
casualidad: cuando se inventó la noción informática de carpeta, se tomó
como modelo la noción matemática de conjunto.
\end{notebox}

\begin{notebox}[Buena práctica, no obligación lógica]
Aunque en teoría un conjunto puede mezclar tipos de objetos distintos
(números, puntos, rectas, etc.), en la práctica se considera \textbf{mala
práctica} hacerlo. El curso trabajará con conjuntos ``tipados'': conjuntos
de números, conjuntos de puntos, etc., sin mezclarlos.
\end{notebox}

\section{Conjuntos: pertenencia e igualdad}

\subsection{Relación de pertenencia y diagramas de Venn}

La noción de conjunto es \textbf{primitiva}: no se define, se observa a
través de una relación, la \textbf{pertenencia}:
\[
x \in A
\]
que se lee ``$x$ es un elemento de $A$'', y su negación
\[
x \notin A.
\]

\begin{notebox}[Asimetría de la notación]
A la derecha del símbolo $\in$ siempre debe haber un \textbf{conjunto}; a
la izquierda, $x$ puede ser \textbf{cualquier objeto matemático} (número,
punto, recta, matriz, función, o incluso otro conjunto --- así como una
carpeta puede contener otras carpetas).
\end{notebox}

Los conjuntos se representan con \textbf{diagramas de Venn} (en Francia,
informalmente, ``diagramas patatoidales'' por su forma de papa): un óvalo
que contiene los elementos adentro; si $x \in A$ se dibuja adentro y si
$y \notin A$ se dibuja afuera.

% Clase 1 --- Relación de pertenencia y diagramas de Venn (§3.1).
% El docente: "en general se dibujan los conjuntos así, con los elementos
% adentro... si A es mi conjunto, tenemos que x pertenece a A, mientras que
% y no pertenece a A". Diagrama de Venn simple: un óvalo (conjunto A) con un
% elemento adentro y uno afuera.
\begin{center}
\begin{tikzpicture}[scale=1]
  % óvalo del conjunto A
  \draw[figset] (0,0) ellipse (2.2 and 1.4);
  \node[figlbl] at (0,1.75) {$A$};

  % elemento x, adentro
  \node[figptoY, minimum size=6pt] (px) at (-0.6,0.25) {};
  \node[figlbl, below=2pt] at (px) {$x$};
  \node[figlbl] at (-0.6,-1.75) {$x \in A$};

  % elemento y, afuera
  \node[figptoZ, minimum size=6pt] (py) at (3.4,0.6) {};
  \node[figlbl, above=2pt] at (py) {$y$};
  \node[figlbl] at (3.4,-1.75) {$y \notin A$};
\end{tikzpicture}\\[0.5em]
{\small\itshape Diagrama de Venn: el conjunto $A$ se dibuja como un óvalo; sus
elementos van adentro ($x$) y los objetos que no pertenecen a $A$ quedan
afuera ($y$).}
\end{center}


\subsection{Igualdad de conjuntos}

La igualdad entre dos conjuntos se \textbf{define} mediante la
pertenencia:
\[
A = B \;:\Longleftrightarrow\; \forall x,\; (x \in A \Longleftrightarrow x \in B)
\]

Es decir, dos conjuntos son iguales si y solo si \textbf{tienen exactamente
los mismos elementos}, sin importar el orden en que se listen ni cuántas
veces se repita cada uno.

\begin{notebox}[Sobre la notación de definición]
El docente distingue la \textbf{igualdad} (relación entre dos \emph{objetos}
matemáticos, que dice que son idénticos) de la \textbf{equivalencia
definicional} (relación entre dos \emph{frases} o enunciados, que dice que
tienen el mismo significado). Como $A = B$ es en sí mismo un enunciado (no
un objeto), no tiene sentido escribir una ``igualdad de igualdades''; por
eso se usa un símbolo distinto ($:\Leftrightarrow$) para introducir la
definición. Quien lo prefiera puede reemplazarlo por un $\Leftrightarrow$
ordinario ``de definición'' --- el significado es el mismo.
\end{notebox}

\begin{notebox}[Cuantificador $\forall x$ sin relativizar]
En esta definición, ``para todo $x$'' cuantifica sobre \textbf{todo el
universo matemático}, no sobre un conjunto particular --- es, en palabras
del docente, ``una cuantificación terrible''.
\end{notebox}

Otra diferencia con las carpetas informáticas: dos carpetas con el mismo
contenido pero distinto nombre son objetos distintos en una computadora;
en teoría de conjuntos, como los conjuntos son anónimos, \textbf{dos
conjuntos con el mismo contenido son el mismo conjunto}.

\section{Inclusión}

\subsection{Definición y cuantificadores relativizados}

La \textbf{inclusión} (subconjunto) se define:
\[
A \subset B \;:\Longleftrightarrow\; \forall x,\; (x \in A \Longrightarrow x \in B)
\]

Nótese que solo cambia un símbolo ($\Rightarrow$ en vez de
$\Leftrightarrow$) respecto de la definición de igualdad. Intuitivamente,
$A \subset B$ dice que todo elemento de $A$ es también elemento de $B$
(``$A$ está incluido en $B$''), en analogía con el orden $\le$ sobre
números.

Esta definición motiva la introducción de los \textbf{cuantificadores
relativizados}, abreviaturas de uso constante en el curso:
\[
\forall x \in A,\; (\ldots) \quad:\Longleftrightarrow\quad \forall x,\; (x \in A \Rightarrow (\ldots))
\]
\[
\exists x \in A,\; (\ldots) \quad:\Longleftrightarrow\quad \exists x,\; (x \in A \;\wedge\; (\ldots))
\]

\begin{notebox}[Cuantificación relativizada vs. no relativizada]
Diferencia entre cuantificador ``no relativizado'' (sobre todo el
universo) y ``relativizado'' (restringido a los elementos de un conjunto
$A$): el segundo es abreviatura del primero combinado con $\Rightarrow$
(para ``para todo'') o $\wedge$ (para ``existe''). El docente insiste en
el símbolo universal $\wedge$ para ``y'' (independiente del idioma, a
diferencia de ``and'', ``et'', ``y''), y en que ``existe'' en matemática
significa \textbf{existe al menos uno}, no ``existen todos''.
\end{notebox}

\subsection{Propiedades de la inclusión}

Para todos los conjuntos $A, B, C$:

\begin{enumerate}
  \item \textbf{Reflexividad}: $A \subset A$ (todo conjunto se incluye a
    sí mismo).
  \item \textbf{Transitividad}: si $A \subset B$ y $B \subset C$,
    entonces $A \subset C$.
  \item \textbf{Antisimetría}: si $A \subset B$ y $B \subset A$, entonces
    $A = B$.
\end{enumerate}

La propiedad (3) es el contenido del

\begin{keybox}[Método estándar para probar igualdad de conjuntos]
\[
\boxed{A = B \iff (A \subset B \;\wedge\; B \subset A)}
\]
Para probar $A = B$ se hacen \textbf{dos demostraciones independientes},
primero $A \subset B$ y luego $B \subset A$, y de ambas se concluye la
igualdad. Este método se usará repetidamente a lo largo del curso.
\end{keybox}

\begin{notebox}[Analogía con el orden numérico]
La inclusión comparte estructura con la relación $\le$ sobre los números:
reflexiva, transitiva y antisimétrica son exactamente los axiomas de un
\textbf{orden (parcial)}.
\end{notebox}

\section{Los conjuntos de base: N, Z, Q, R}

El curso trabajará principalmente con cuatro conjuntos numéricos de base,
relacionados por inclusiones:
\[
\boxed{\mathbb{N} \subset \mathbb{Z} \subset \mathbb{Q} \subset \mathbb{R}}
\]

\begin{table}[H]
\centering
\begin{tabularx}{\textwidth}{@{}llXX@{}}
\toprule
Conjunto & Notación & Descripción informal & Definición rigurosa \\
\midrule
Naturales & $\mathbb{N}$ & $0, 1, 2, 3, \ldots$ (en este curso empiezan en $0$) & Pendiente (cap. 2) \\
Enteros relativos & $\mathbb{Z}$ & $\ldots, -2, -1, 0, 1, 2, \ldots$ & Pendiente (cap. 2) \\
Racionales & $\mathbb{Q}$ & Fracciones & $\mathbb{Q} = \{\, \tfrac{n}{p} : n, p \in \mathbb{Z},\ p \neq 0 \,\}$ \\
Reales & $\mathbb{R}$ & --- & Se define recién en el capítulo 2 \\
\bottomrule
\end{tabularx}
\end{table}

\begin{warnbox}[Notaciones ``flojas'' con puntos suspensivos]
Escribir $\mathbb{N} = \{0, 1, 2, 3, \ldots\}$ no es una definición
matemática precisa porque los puntos suspensivos ocultan una cantidad
infinita de elementos que nunca se terminarían de listar (el docente
ilustra la idea con la imagen, en broma, de escribir naturales hasta
destruir las paredes del salón y dar la vuelta al planeta, sin terminar
nunca). $\mathbb{Z}$ es ``una broma doble'' porque tiene infinitos en dos
direcciones, y $\mathbb{Q}$ es peor todavía porque además hay una cantidad
infinita de racionales entre dos racionales cualesquiera. Por eso
$\mathbb{Q}$ sí admite una definición por comprensión rigurosa (fracciones
de enteros con denominador no nulo), mientras que $\mathbb{N}$ y
$\mathbb{Z}$ requieren una construcción más cuidadosa que se pospone al
capítulo 2.
\end{warnbox}

Sobre $\mathbb{R}$, el docente adelanta que la diferencia con $\mathbb{Q}$
no es solo agregar ``algunos números para completar'' (idea que trae el
estudiante típico del liceo, y que es \textbf{falsa}): hay
\textbf{infinitamente más} números en $\mathbb{R}$ que en $\mathbb{Q}$.
Ejemplos de reales no racionales: $\sqrt{2}$, $\sqrt{3}$, $-\sqrt{2}$,
$\pi$, $e$.

\begin{notebox}[Nota histórica]
El docente sitúa el nacimiento del análisis matemático (la disciplina de
Cálculo 1) en los trabajos de \textbf{Leibniz} hacia fines del siglo XVII
(contemporáneo de Bach y Händel). En esa época aún no se sabía qué era
rigurosamente un número real, lo que generó ``cuentas flojas'' y
paradojas; la definición precisa y moderna de $\mathbb{R}$ es la tarea del
capítulo 2 (dos semanas de curso), a partir de la cual $\mathbb{N}$,
$\mathbb{Z}$ y $\mathbb{Q}$ se recuperan como subconjuntos particulares.
\end{notebox}

\section{Construcción de conjuntos por extensión}

Sin ejemplos de conjuntos concretos aún, el docente introduce dos métodos
generales para \textbf{construir} conjuntos, empezando por la
\textbf{construcción por extensión}: dar la lista finita de todos sus
elementos.

Dados objetos cualesquiera $a_1, \ldots, a_n$, se define
\[
\{a_1, \ldots, a_n\}
\]
por la propiedad
\[
\forall x,\; \big(x \in \{a_1,\ldots,a_n\} \iff x = a_1 \lor x = a_2 \lor \cdots \lor x = a_n\big).
\]

\begin{notebox}[La disyunción $\lor$ es inclusiva por defecto]
El símbolo $\lor$ (disyunción) es, en la convención matemática por
defecto, \textbf{inclusivo}: ``$A$ o $B$'' admite que ambas se cumplan a
la vez. Es el símbolo ``dual'' de $\wedge$: gráficamente, $\wedge$
invertido da $\lor$. Cuando se necesita disyunción exclusiva, se escribe
la palabra completa ``exclusivo'', porque no hay símbolo estándar para
ella.
\end{notebox}

\textbf{Defecto de este método}: solo permite construir \textbf{conjuntos
finitos} (hay que poder listar todos los elementos), y en principio no
impone ninguna restricción sobre el tipo de los $a_i$ (podrían mezclarse
números, puntos, conjuntos, etc. --- aunque en la práctica no se hace, por
las razones de ``buena práctica'' ya mencionadas en la Sección 2).

Como un conjunto no puede contener dos veces el mismo elemento, listar un
objeto repetido en la construcción no cambia el resultado: $\{a, a\} =
\{a\}$.

El docente distingue tres \textbf{casos particulares} según la cantidad
$n$ de elementos listados.

\subsection{Caso $n=0$: el conjunto vacío}

\begin{keybox}[Definición del conjunto vacío]
\[
\boxed{\varnothing := \{\,\} \quad\text{tal que}\quad \forall x,\; x \notin \varnothing}
\]
\end{keybox}

El vacío es, por definición, el conjunto construido por extensión a
partir de una \textbf{lista vacía}. Corresponde a una \textbf{disyunción
vacía}, que por convención (análoga a la de una sumatoria vacía igual a
$0$) se toma como \textbf{falsa} --- el elemento neutro de la disyunción.

\subsection{Caso $n=1$: el conjunto unitario}

$\{a\}$ se llama \textbf{conjunto unitario}.

\begin{warnbox}[$a$ y $\{a\}$ son objetos distintos]
Frecuente fuente de confusión en las pruebas: $a$ y $\{a\}$ son objetos
\textbf{distintos} --- como la diferencia entre un archivo $a$ y una
carpeta que contiene únicamente el archivo $a$ (o entre una manzana y una
bolsa que contiene una manzana). Esto vale \textbf{incluso cuando $a$ es
en sí mismo un conjunto}: el ejemplo canónico es
\[
\varnothing \neq \{\varnothing\},
\]
el conjunto vacío frente al conjunto unitario que contiene al vacío. El
primero no tiene elementos; el segundo tiene exactamente un elemento (que
resulta ser el conjunto vacío). Distinguir estos ``niveles'' es, según el
docente, una de las sutilezas más comunes de confundir por quien empieza
a estudiar teoría de conjuntos.
\end{warnbox}

\subsection{Caso $n=2$: el par no ordenado}

$\{a, b\}$ se llama \textbf{par no ordenado}, porque
\[
\{a, b\} = \{b, a\}
\]
dado que ambos conjuntos tienen los mismos elementos (el orden en que se
listan los elementos de un conjunto es irrelevante, igual que el orden de
los archivos dentro de una carpeta). Como caso límite,
\[
\{a, a\} = \{a\},
\]
es decir, listar el mismo objeto dos veces produce el conjunto unitario,
no un ``conjunto con dos copias''.

\section{Construcción de conjuntos por comprensión}

El método por extensión no sirve para conjuntos infinitos, que son los de
mayor interés en matemática. El segundo método, \textbf{construcción por
comprensión}, resuelve esto: dado un conjunto $A$ ya formado y un
\textbf{enunciado} (frase, no objeto matemático) $P(x)$ que depende de una
variable $x$, se define
\[
\{\, x \in A : P(x) \,\}
\]
por la propiedad
\[
\forall y,\; \big(y \in \{x \in A : P(x)\} \iff (y \in A \;\wedge\; P(y))\big).
\]

\begin{notebox}[El símbolo $x$ está ligado]
El símbolo $x$ dentro de las llaves está \textbf{ligado}: solo tiene
sentido dentro de la notación $\{x \in A : P(x)\}$, igual que la variable
ligada por $\forall$ o $\exists$, o como la $x$ en la definición de una
función ``$x \mapsto 3x$''. El conjunto resultante \textbf{no depende de
$x$}, solo de $A$ y de la frase $P$; se puede sustituir $x$ por cualquier
nombre sin cambiar el conjunto (y se puede evaluar la pertenencia de
cualquier objeto concreto, no solo de una variable, p.\,ej.\ $3 \in \{x
\in A: P(x)\} \iff 3 \in A \wedge P(3)$).
\end{notebox}

\textbf{Interpretación geométrica}: la propiedad $P$ traza una
\textbf{frontera} en el universo matemático, separando los objetos que la
cumplen de los que no. La construcción por comprensión ``recorta'', dentro
de $A$, exactamente los elementos que caen del lado de la frontera donde
$P$ es verdadera.

% Clase 1 --- Construcción de conjuntos por comprensión (§7).
% El docente: "supongo que yo tengo un conjunto A... una propiedad me va a
% dividir el universo en dos partes: los objetos que cumplen la propiedad y
% los que no... ¿cuáles serían los elementos que están dentro de A pero que
% no cumplen P de x? Un elemento acá, por ejemplo Y que cumple P(x) y este Z
% que no cumple P(x)". El universo matemático se parte en dos mitades según
% P(x); el conjunto A se solapa con ambas; solo la intersección con la mitad
% "cumple P" es el conjunto {x en A : P(x)}.
\begin{center}
\begin{tikzpicture}[scale=1]
  % ---- universo matemático ----
  \draw[figuniverso] (-4,-2.5) rectangle (4,2.5);
  \node[figlblaux] at (-3.3,2.2) {universo};

  % ---- conjunto A (base, antes de resaltar la intersección) ----
  \draw[figset] (1,0) ellipse (2 and 1.3);
  \node[figlbl] at (1,1.7) {$A$};

  % ---- frontera P(x): divide el universo en dos mitades ----
  \draw[figaux] (0,-2.5) -- (0,2.5);
  \node[figlblaux, above] at (0,2.5) {$P(x)$};
  \node[figlblaux] at (2.4,-2.15) {cumple $P$};
  \node[figlblaux] at (-2.4,-2.15) {no cumple $P$};

  % ---- resaltar la región intersección: A \cap "cumple P" ----
  \begin{scope}
    \clip (0,-2.5) rectangle (4,2.5);
    \clip (1,0) ellipse (2 and 1.3);
    \fill[figamber!25] (-4,-2.5) rectangle (4,2.5);
    \draw[figsetC] (1,0) ellipse (2 and 1.3);
    \draw[figsetC] (0,-2.5) -- (0,2.5);
  \end{scope}

  % ---- elemento y: en A y cumple P (entra al conjunto por comprensión) ----
  \node[figptoY, minimum size=6pt] (py) at (1.5,-0.3) {};
  \node[figlbl, right=2pt] at (py) {$y$};

  % ---- elemento z: en A pero no cumple P (queda fuera) ----
  \node[figptoZ, minimum size=6pt] (pz) at (-0.5,0.35) {};
  \node[figlbl, above=2pt] at (pz) {$z$};
\end{tikzpicture}\\[0.5em]
{\small\itshape La propiedad $P$ traza una frontera que parte el universo en
dos mitades; el conjunto $\{x \in A : P(x)\}$ es la región resaltada
(intersección de $A$ con la mitad ``cumple $P$''). Tanto $y$ como $z$
pertenecen a $A$, pero solo $y$ cumple $P$ y entra al nuevo conjunto; $z$ no
cumple $P$ y por eso queda afuera, aunque también pertenezca a $A$.}
\end{center}


\begin{notebox}[Los dos ingredientes son indispensables]
No se puede prescindir de $A$ y tomar ``todos los objetos del universo que
cumplen $P$'', porque en general esa colección sería demasiado grande. Si
$A$ es finito, el conjunto resultante también lo es; si $A$ es infinito,
el resultado puede ser finito o infinito según la propiedad $P$.
\end{notebox}

\textbf{Ejemplos trabajados en clase} (todos con $A = \mathbb{N}$):

\begin{table}[H]
\centering
\begin{tabularx}{\textwidth}{@{}lXX@{}}
\toprule
Propiedad $P(n)$ & Conjunto resultante & ¿Finito o infinito? \\
\midrule
$n$ es par & $\{0, 2, 4, 6, \ldots\}$ & Infinito \\
$n$ es impar & $\{1, 3, 5, 7, \ldots\}$ & Infinito (complementario del anterior) \\
$n$ es primo & $\{2, 3, 5, 7, 11, 13, \ldots\}$ & Infinito (estructura fina aún conjetural) \\
$n$ es primo y par & $\{2\}$ & Finito (un único elemento) \\
$n$ es par y $n$ es impar & $\varnothing$ & Vacío \\
\bottomrule
\end{tabularx}
\end{table}

\begin{notebox}[El vacío también se construye por comprensión]
El último ejemplo muestra que \textbf{el conjunto vacío también puede
construirse por comprensión}: basta una propiedad que ningún objeto del
universo satisface. También se señaló, a partir de una pregunta de la
clase, que \textbf{negar la propiedad} $P$ (usar $\neg P$ en vez de $P$)
produce el \textbf{conjunto complementario} dentro de $A$ --- como en el
par par/impar de arriba, que son complementarios entre sí dentro de
$\mathbb{N}$.
\end{notebox}

La clase termina en este punto (``vamos a continuar con la comprensión la
próxima vez''), habiendo cubierto pertenencia, igualdad, inclusión, los
conjuntos numéricos de base y los dos métodos de construcción de conjuntos
(extensión y comprensión), como preparación para retomar funciones y
avanzar hacia la definición rigurosa de $\mathbb{R}$ en el capítulo 2.

\emph{Continúa en la Clase 2 con más ejemplos de construcción por
comprensión y, previsiblemente, la noción de función.}

\end{document}
